锂离子电池凭借高能量密度、长循环寿命和低自放电率等优势，已广泛应用于电动汽车、便携式电子设备和规模化储能等领域\cite{ref1,ref2,ref3,ref4}，是支撑新能源交通、智能终端和能源存储行业发展的重要储能技术。尽管如此，电池在长期运行中仍存在不可忽视的安全风险。充放电循环与复杂工况会引发固体电解质界面（solid electrolyte interphase, SEI）膜增厚、活性锂损失和电极结构衰退等老化现象\cite{ref5,ref6}，进而导致容量衰减和性能下降，严重时将诱发短路、热失控等故障\cite{ref7,ref8}。因此，精确估计锂离子电池的健康状态对于保障电池的安全性、可靠性和运行性能至关重要。

容量衰减是电池老化最直观的特征，健康状态（state of health, SOH）通常采用容量保持率定义\cite{ref13,ref14}。但电池最大可用容量的精确测算需要完整或近似完整的充放电循环，难以满足在线监测需求\cite{ref10,ref11}。因此，从可观测运行信号中间接估计SOH已成为在线健康监测的关键途径。然而，复杂工况增加了退化信息稳定提取的难度，而电池管理系统（battery management system, BMS）严格的计算与存储限制又制约了模型复杂度，使在线SOH估计面临退化信息提取与模型计算效率的双重挑战\cite{ref31}。

\subsection{文献综述}

现有锂离子电池健康状态（SOH）估计方法分为模型驱动和数据驱动两类\cite{ref10,ref11,ref12}。前者包括电化学模型和等效电路模型，后者包括传统机器学习和深度学习方法。

模型驱动方法使用数学方程或等效电路模拟电池内部的电化学反应机制\cite{ref19}。电化学模型利用微分方程描述电池老化过程，其状态变量和模型参数具有明确的物理含义。例如，Li等人\cite{ref21}将SEI膜生长和电极颗粒开裂纳入单粒子模型，用于SOH估计。这类模型虽具有较好的物理可解释性和估计精度，但复杂方程求解计算开销较大，且颗粒半径、扩散系数等微观参数获取难度大，限制了其在SOH在线估计中的应用\cite{ref20,ref22,ref23}。

相比之下，等效电路模型（ECM）是一类应用广泛的简化电学模型，使用电阻和电容等元件模拟电池充放电动态\cite{ref19}。该类模型计算复杂度较低，参数辨识也相对容易。例如，Chen等人\cite{ref24}采用递推最小二乘法辨识Thevenin模型参数，并根据欧姆内阻与容量衰减的关系估计SOH。然而，ECM无法全面刻画电池内部状态的演变规律，估计精度高度依赖于电路结构与模型参数。同时，模型参数对温度、充放电倍率等工况变化较为敏感，使ECM难以在整个老化周期内保持良好的估计鲁棒性\cite{ref19}。

相较于模型驱动方法，数据驱动方法利用历史运行数据及其构造的健康指标，学习外部可观测信息与SOH之间的映射关系\cite{ref11,ref14}，降低了对复杂物理模型和微观参数的依赖。早期研究主要采用支持向量机、随机森林和高斯过程回归等传统机器学习模型\cite{ref25}。Fei等人\cite{ref26}从前100个循环的电压、电流、温度、内阻和容量数据中构造42个特征，并通过特征筛选比较了六种机器学习模型的早期寿命预测结果。Li等人\cite{ref27}采用随机森林，根据局部充电数据估计电池容量。传统机器学习模型结构相对简单，但通常依赖人工构造的健康指标，其估计性能取决于输入特征能否稳定反映电池退化。

常用健康指标主要从电池循环过程中的电压、电流、温度、时间和内阻等运行数据及其衍生曲线中构造，不同指标在可获取性、工况稳定性和预处理要求上存在差异\cite{ref47,ref48,ref49,ref50}。部分内阻指标需要进行阻抗谱测试，该过程耗时且依赖专用仪器\cite{ref64}。温度指标包含电池退化信息，但容易受到环境温度和充放电倍率的影响\cite{ref64}。增量容量曲线反映充电电压平台的变化，微分温度--电压曲线和微分温度--容量曲线则显示充电过程中的热响应变化\cite{ref50,ref64,ref65}。例如，Wen等人\cite{ref50}提取了增量容量曲线的峰值、峰位和峰形斜率，Dai等人\cite{ref48}从电压、电流、温度、增量容量和微分温度--电压曲线中构造了多类统计特征。计算这类微分曲线需要对采样数据进行求导，但微分运算会放大噪声和相邻采样点的波动，因此通常需要先对采样数据进行平滑处理\cite{ref28,ref50}。

相比之下，时间类健康指标无需微分处理，提取过程更为简单。Richardson等人\cite{ref28}从恒流充电阶段的多个电压区间中提取时间特征。恒流充电时间（CCCT）是其中一种常用特征，只需计算选定电压区间两端的时间差。例如，Lin等人\cite{ref63}采用恒流充电时间估计电池SOH。电压区间改变后，该指标与容量之间的相关性也会改变。Tian等人\cite{ref51}通过搜索充电电压区间提高时间特征与容量之间的相关性，Li等人\cite{ref31}则逐步缩小充电电压区间，并利用皮尔逊相关系数筛选候选区间。这些方法通常根据选定电池的相关结果确定电压区间；若窗口由单节电池确定，其在同组其他电池上的相关性稳定性仍需检验。因此，区间选择需综合多节电池的相关性结果。

在预测模型方面，卷积神经网络（CNN）、循环神经网络（RNN）和Transformer等深度学习方法已被用于电池SOH估计\cite{ref14}。CNN利用卷积核提取局部特征。例如，Qian等人\cite{ref29}采用一维CNN根据充电曲线片段估计电池容量，而Lee等人\cite{ref30}则将跨循环容量衰减序列重构为矩阵，并采用二维CNN估计SOH。然而，传统CNN并不擅长处理长时间序列，因为每层卷积通常只在有限的局部感受野内聚合信息，相距较远的特征需要经过多层网络才能融合\cite{ref31}。RNN则通过递推更新隐藏状态，使历史信息能够沿时间步持续传递；LSTM和GRU进一步利用门控机制选择性保留和更新这些信息，从而捕捉长期依赖和时序动态\cite{ref14}。尽管如此，循环状态仍需逐步更新，各时间步难以并行计算，长序列训练和推理的吞吐量因而受到限制\cite{ref31}。针对上述局限，Vaswani等人\cite{ref34}提出了Transformer。该模型避免循环，转而利用自注意力捕捉全局依赖关系，并支持序列位置的并行计算。例如，Park和Kim\cite{ref38}采用物理先验引导的Transformer估计电池长期SOH，Chen等人\cite{ref37}则将视觉Transformer用于SOH估计。

为兼顾局部特征提取和长期依赖建模，研究者进一步将不同网络结构进行组合。CNN-LSTM、CNN-BiLSTM和CNN-GRU结合了卷积与循环结构\cite{ref33,ref41,ref42}，BiGRU-Transformer、CNN-Transformer和卷积Transformer则进一步引入了自注意力机制\cite{ref32,ref35,ref36}。例如，Xu等人\cite{ref41}采用CNN提取局部特征，再利用LSTM建模时间依赖；Gu等人\cite{ref35}将CNN与Transformer结合，以融合局部特征和长程依赖。这些混合结构补充了单一模型的能力，但额外模块也增加了参数量和计算开销\cite{ref31}。同时，标准自注意力需要计算所有序列位置之间的两两关系，其时间和存储开销随序列长度呈二次增长\cite{ref39,ref40}。上述资源开销增加了模型在资源受限BMS中的在线应用难度\cite{ref31}。

\subsection{挑战分析与贡献}

据此，当前数据驱动锂离子电池SOH估计面临的主要挑战可归纳为以下三点。

（1）\textbf{健康指标构建缺乏跨电池稳定性。} 从多传感器收集的原始数据包含多样化的退化相关特征；然而，许多现有方法仍依赖单一或有限的数据源。此外，CCCT通常从预设电压区间中提取，其窗口调整和候选健康指标筛选主要依据选定电池上的PCC或SCC。单独使用一种相关系数只能考察线性关系或单调关系。只比较候选指标与SOH之间的相关性，也不能识别指标之间的重复信息。不同电池上的相关性结果还可能发生变化，影响模型输入的稳定性。

（2）\textbf{单一结构难以兼顾短期变化和长期趋势。} 随着循环推进，电池容量整体呈下降趋势，并伴随相邻循环的波动和局部容量恢复。CNN主要提取相邻循环的局部变化，RNN和Transformer主要捕捉跨循环的长期变化。模型侧重局部变化时，不同退化阶段之间的联系难以充分建立；模型侧重长期变化时，相邻循环的波动和局部容量恢复容易被弱化。单一结构获得的退化信息因此不够完整，影响SOH估计的准确性。

（3）\textbf{模型精度提升伴随计算开销增加。} 许多现有深度学习模型通过增加网络深度或组合不同网络结构提高SOH估计精度。这些方法增加了模型参数和计算操作。循环结构需要按时间步依次计算，标准自注意力需要计算所有序列位置之间的两两关系，其时间和存储开销随序列长度呈二次增长。这些开销加重了资源受限BMS的运行负担，限制了在线SOH估计模型的部署。

为应对上述挑战，本文提出一种以MS-AgentNet为核心的锂离子电池SOH估计框架。该框架从健康指标构建和轻量模型设计两个层面展开，分别通过组级标定与筛选提高模型输入的跨电池稳定性，并以较低计算开销提取局部变化和长期趋势。本文进一步在四个公开数据集的独立评价电池上检验上述设计。主要贡献如下。

（1）\textbf{提出组级健康指标选择算法。} 该算法以特征开发电池集合上的相关性结果为依据，利用MS-CCCT标定充电时间特征的电压窗口，再通过PCC/SCC双阈值与冗余约束筛选候选指标。统一的筛选规则为各数据集确定相应的健康指标组合，入选指标在独立评价电池上仍保持较强的线性和单调相关性。

（2）\textbf{构建轻量级局部—全局网络MS-AgentNet。} 核心的局部—全局融合注意力用于捕获电池退化序列中的短期变化和长期趋势，并将注意力计算复杂度从二次$O(N^2)$降至线性$O(N)$。模型同时引入大核深度可分离卷积，与局部卷积共同提取不同尺度的退化特征，并减少卷积参数。上述设计共同形成紧凑的模型结构，增强了对电池退化特征的表征能力。

（3）\textbf{开展多数据集综合验证。} 在具有不同材料、容量和充放电协议的多个公开数据集上开展独立电池评价，并进行模块消融、复杂度分析和跨数据集迁移实验，以评估所提方法的估计精度、计算效率和泛化能力。结果表明，所提框架在SOH估计精度和计算效率方面均优于多种对比模型，目标域观测进一步降低了跨数据集迁移误差。
